% !TeX spellcheck = en_US
\documentclass[11pt,a4paper]{article}

\usepackage[T1]{fontenc}
\usepackage[utf8]{inputenc}
\usepackage[margin=2.5cm]{geometry}
\usepackage{booktabs}
\usepackage{amsmath}
\usepackage{amssymb}
\usepackage{siunitx}
\usepackage[hidelinks]{hyperref}
\usepackage{xcolor}
\usepackage{listings}

\lstset{
  basicstyle=\ttfamily\small,
  breaklines=true,
  columns=fullflexible,
  keepspaces=true,
}
\newcommand{\code}[1]{\texttt{#1}}

\title{Refactoring and Performance Optimization of OPENSC2:\\
Modularization, Modernization and Computational Efficiency}
\author{Laura A. M. D'Angelo}
\date{%
	{\small Stellarator Reactor Studies, Max Planck Institute for Plasma Physics, Greifswald (Germany)}\\[2ex]%
	July 2026}


\begin{document}
\maketitle

\begin{abstract}
This document summarizes the restructuring of the OPENSC$^2$ code base (a coupled
thermal-hydraulic-electric solver for superconducting cables) carried out on the
\code{develop} branch. The work comprises three threads: (i) a \emph{modularization}
of the formerly monolithic source tree into physics-oriented packages with
single-responsibility modules; (ii) a \emph{modernization} of the internal data model,
replacing string-keyed dictionaries and Fortran-heritage identifiers with typed
data classes, enumerations and descriptive naming; and (iii) a systematic
\emph{computational-efficiency} campaign covering sparse-matrix assembly, vectorized
finite-element construction, LAPACK-based linear algebra, tabulated fluid properties,
and a new adaptive time-integration layer with second-order time steppers.
The net effect on the W7-X quench benchmark (17\,000 linear elements) is a reduction
of the wall time for the full \SI{25}{s} transient from several hours to
\textbf{19 minutes}, while every individual optimization step was verified to be bit-compatible (relative deviations $\lesssim 10^{-12}$, i.e.\ floating-point summation-order noise) against a frozen solver state, and the physical results reproduce the THEA reference solution.
\end{abstract}

\tableofcontents

%=====================================================================
\section{Starting Point and Goals}
%=====================================================================

The refactoring lives on the \code{develop} branch, which branched off \code{main}
(commit \code{dfeda15}) in June~2026. At that point OPENSC$^2$ consisted of a flat
\code{source\_code/} directory dominated by a few very large modules:
\code{channel.py} (1\,281 lines), \code{conductor.py} (eventually 3\,457 lines,
$\sim$60 methods mixing mesh geometry, topology discovery, electric preprocessing,
heat-transfer physics and diagnostics), \code{simulation.py} (876 lines mixing the
time-stepping loop with output-folder bookkeeping), and
\code{step\_matrix\_construction.py} (1\,473 lines containing the entire
finite-element assembler for all physics at once). State was passed around in
string-keyed dictionaries (\code{dict\_node\_pt}, \code{dict\_Gauss\_pt},
\code{dict\_band}, \code{dict\_Step}, \ldots), physical models were selected by
comparing raw integer flags from the Excel input workbooks, and many identifiers
retained the names of the original Fortran routines (\code{SYSMAT}, \code{ASCALING},
\code{UPWEQT}, \code{gredub}/\code{gbacsb}).

The refactoring pursued three goals, in this order of priority:
\begin{enumerate}
  \item \textbf{Modularization}: one package per physical or infrastructural concern;
        one module per responsibility; orchestration separated from physics.
  \item \textbf{Modernization}: typed data structures, explicit enumerations,
        full descriptive naming (no abbreviations), removal of dead code.
  \item \textbf{Efficiency}: profiling-driven elimination of the dominant costs,
        with a strict non-regression gate at every step.
\end{enumerate}

%=====================================================================
\section{Chronology}
%=====================================================================

The work proceeded in five phases, traceable in the git history of \code{develop}
(Table~\ref{tab:chronology}).

\begin{table}[ht]
  \centering
  \caption{Phases of the refactoring as recorded in the git history of
           \code{develop}.}
  \label{tab:chronology}
  \begin{tabular}{llp{8.6cm}}
    \toprule
    Phase & Commits (period) & Content \\
    \midrule
    0 & \code{3a75265}\,\ldots\,\code{b0768a6} (Jun 2--3) &
        Preparation: \code{pyproject.toml} packaging and package
        \code{\_\_init\_\_} files; Linux/GUI compatibility; newer
        pandas/numpy compatibility; robustness fixes; headless (no-GUI)
        drivers for reproducible runs. \\
    1 & \code{9346a68} (Jun 16) &
        Channel decomposition: the 1\,281-line \code{channel.py} split into a
        package of friction-factor and heat-transfer \emph{model classes} with
        factories (strategy pattern). \\
    2 & \code{4dfb626} ``part 1'' (Jul 3) &
        The main modularization: 88 files, $+13\,447/-10\,458$ lines ---
        package split (Sec.~\ref{sec:modularization}), typed data model
        (Sec.~\ref{sec:modernization}), runtime migration to the point where the
        modularized code executes end-to-end. \\
    3 & \code{fe68476}, \code{fba79e3} ``parts 2--3'' (Jul 3) &
        Performance campaign: factor 6.7 (electric module, LAPACK banded solver),
        then another factor $\sim$2.8 with full assembler vectorization, plus
        de-dictification and descriptive renaming
        (Sec.~\ref{sec:efficiency}). \\
    4 & \code{93fa4d8}, \code{9e0ec1e} ``part 4'' (Jul 4) &
        Tabulated fluid properties; new time integrators (Galerkin $\theta=2/3$,
        variable-step BDF2) and adaptive time stepping with local-truncation-error
        control (Sec.~\ref{sec:timestepping}). \\
    \bottomrule
  \end{tabular}
\end{table}

%=====================================================================
\section{Modularization Strategy}
\label{sec:modularization}
%=====================================================================

\subsection{Package structure}

The flat layout was reorganized into physics-oriented packages
(Table~\ref{tab:packages}). The guiding rule: a module belongs to the package that
owns its \emph{physics}, not to the module that happens to call it.

\begin{table}[ht]
  \centering
  \caption{Package layout of the modularized \code{source\_code/} tree.}
  \label{tab:packages}
  \begin{tabular}{ll}
    \toprule
    Package & Responsibility \\
    \midrule
    \code{components/}          & Component classes (fluid, strands, stacks, jackets) and factories \\
    \code{conductor/}           & Conductor orchestration, mesh, topology, input loading/validation \\
    \code{thermal/}             & Energy equation, conduction, heat sources, radiation, HTC models \\
    \code{hydraulics/}          & Momentum equation, coolant state, channel/friction models \\
    \code{electromagnetics/}    & Electric network model: topology, resistance, conductance, solver \\
    \code{physical\_fields/}    & Typed field containers (nodal/Gauss fields, time evolutions) \\
    \code{interfaces/}          & External-library interfaces (CoolProp fluid properties) \\
    \code{properties\_of\_materials/} & Material property correlations \\
    \code{utility\_functions/}  & Generic FEM assembly, transient driver functions, I/O, paths \\
    \bottomrule
  \end{tabular}
\end{table}

\subsection{Physical models as classes: the channel decomposition}

The first structural refactoring (phase 1) established the pattern for
\emph{selectable physical models}: the monolithic \code{channel.py} --- friction
factors and heat-transfer correlations interleaved with channel bookkeeping in
1\,281 lines --- was decomposed into dedicated model classes (laminar, turbulent and
total friction correlations; Nusselt correlations) instantiated through factory
classes (\code{FrictionFactorFactory}, \code{HeatTransferFactory}) that translate the
workbook flag once into a model object (strategy pattern). Adding a correlation is now
a local change: one model class plus one factory branch --- exercised in practice when
the bare Dittus--Boelter model and the Darcy--Forchheimer porous-medium friction were
wired in for the W7-X benchmark. These modules were later relocated into
\code{thermal/} (heat transfer) and \code{hydraulics/} (friction) when those packages
were created.

\subsection{Splitting the monoliths}

The decomposition followed a single, uniform calling convention established early and
then reused everywhere: physics is implemented as \emph{free functions taking the
\code{conductor} (and where needed \code{simulation}) object as explicit first
argument}; the \code{Conductor} class keeps only thin one-line wrapper methods where
external callers (GUI, headless drivers) require a stable public API. Highlights:

\begin{itemize}
  \item \textbf{Thermal consolidation}: the energy equation was extracted from the
        monolithic assembler into \code{thermal/} (\code{energy\_equation.py},
        \code{conduction.py}, \code{thermal\_transport.py}, \code{heat\_sources.py},
        \code{radiation.py}, \code{temperature\_field.py}, Nusselt/HTC models), with the
        velocity/pressure rows split into the counterpart
        \code{hydraulics/momentum\_equation.py}. Where one function had computed all
        equation rows while sharing intermediate terms, the split preserved the required
        \emph{momentum-before-energy call order} on the shared matrix, documented in both
        docstrings. \code{step\_matrix\_construction.py} shrank from 1\,473 to
        $\sim$500 lines and now contains only physics-agnostic FEM assembly.
  \item \textbf{Conductor/simulation split}: \code{conductor.py} went from 3\,457 to
        1\,241 lines by extracting the $\sim$618-line heat-transfer-coefficient method
        into four per-interface-kind functions (\code{thermal/htc\_evaluation.py}), the
        mutually recursive component-topology discovery into
        \code{conductor/conductor\_topology.py}, and the mesh construction into the
        existing \code{conductor/conductor\_mesh.py}. \code{simulation.py} went from 876
        to 484 lines by moving all output-folder bookkeeping into
        \code{utility\_functions/simulation\_paths.py}; the class retains only the
        simulation loop --- its actual single responsibility.
  \item \textbf{Electric module}: the electric network solver had already been
        decomposed into \code{electromagnetics/} submodules (circuit topology,
        resistance, conductance, system assembly, boundary conditions, solver); the
        14 corresponding \code{Conductor} methods are one-line delegations and served
        as the architectural template for the rest of the split.
\end{itemize}

\subsection{Managing import cycles}

Two systematic measures keep the package graph acyclic in practice:
\code{electromagnetics/\_\_init\_\_.py} uses PEP~562 lazy re-exports
(module-level \code{\_\_getattr\_\_}) instead of eager imports, and modules inside
\code{conductor/} or \code{thermal/} that need a utility module which itself imports
\code{Conductor} defer that import into the function body. Both patterns are documented
in the code as the house rule for future additions.

%=====================================================================
\section{Modernization}
\label{sec:modernization}
%=====================================================================

\subsection{Packaging and platform compatibility}

The preparation phase turned the source tree into a proper installable package
(\code{pyproject.toml}, package \code{\_\_init\_\_} files), fixed Linux
compatibility of the GUI (including deprecated file-dialog calls), restored
compatibility with current pandas/numpy releases (positional-indexing semantics,
\code{float()} on length-one arrays), and added \emph{headless drivers} so that every
test case runs reproducibly without the GUI --- the precondition for all later
automated verification.

\subsection{Typed data structures instead of string-keyed dictionaries}

\begin{itemize}
  \item \textbf{Field containers}: the pervasive \code{dict\_node\_pt}/%
        \code{dict\_Gauss\_pt} dictionaries ($\sim$600 call sites across all physics)
        were hard-renamed to \code{node\_fields}/\code{gauss\_fields}, instances of
        \code{physical\_fields.FieldContainer} with typed attribute access
        (\code{fields.temperature} instead of \code{dict["temperature"]}), field-unit
        metadata, and time-evolution recording via declared per-component field lists.
  \item \textbf{Solver structures} (\code{conductor/solver\_structures.py}): the former
        \code{dict\_N\_equation}, \code{dict\_band}, \code{dict\_Step} and
        \code{dict\_norm} became the data classes \code{EquationCounts},
        \code{BandStructure}, \code{TimeIntegrationState} and \code{SolutionNorms}.
        The assembler bundles \code{SystemMatrices}, \code{GaussPointMatrices} and
        \code{ElementMatrices} replaced the \code{MASMAT}/\code{FLXMAT}/\ldots{}
        dictionaries; \code{SystemMatrices} implements a custom \code{\_\_iter\_\_}
        yielding the array attributes themselves so that in-place assembly mutation is
        preserved (a \code{dataclasses.astuple}-based implementation would silently
        deep-copy).
  \item \textbf{Coupling data}: the never-populated \code{dict\_df\_coupling} was
        replaced by \code{conductor.coupling}, a \code{ConductorCoupling} of
        \code{CouplingMatrix} objects supporting name-or-index access with
        upper-triangle folding.
\end{itemize}

\subsection{Enumerations for input-workbook flags}

All raw integer/string flags from the Excel workbooks are now translated once, at load
time, into explicit \code{IntEnum}s and compared symbolically thereafter:
\code{MethodFlag} (time integrators), \code{HTC\_Choice}, \code{HeatExcitation},
\code{ContactPerimeterFlag}, \code{HeatTransferModelType}, \code{HydraulicBC},
\code{BFieldDefinitionType}, \code{ElectricConductanceMode}, among others. Method
families are expressed as tuples (\code{THETA\_FAMILY\_METHODS},
\code{ONE\_STEP\_METHODS}) so that array-layout branches never enumerate individual
members.

\subsection{Naming and dead-code policy}

Every new or renamed identifier uses full descriptive words
(\code{evaluate\_transport\_coefficients}, not \code{eval\_transp\_coeff});
Fortran-heritage names were retired (\code{SYSMAT} $\to$ \code{system\_matrix},
\code{Known} $\to$ \code{known\_term\_vector},
\code{ASCALING} $\to$ \code{row\_scaling\_factors},
\code{UPWEQT} $\to$ \code{upwind\_weights}, \code{EIGTIM} $\to$ 
\code{time\_accuracy\_eigenvalue}, \ldots). Dead code was deleted rather than carried
along, e.g.\ the hard-coded profiling scaffold in \code{simulation.py} and the
hand-written banded LU pair \code{gredub}/\code{gbacsb}.

\subsection{Robustness}

Fail-fast guards with full state dumps (\code{.npz}) were added around the thermal
solve: non-finite solution components, reconstructed fluid fields leaving the physical
domain ($p \le 0$, $T \le 0$), and invalid inputs to the fluid-property interface all
raise immediately with the offending node, time and a reproducible state file ---
previously such states propagated silently into misleading downstream crashes.
A related robustness fix nudges temperatures off the CoolProp critical-temperature
singularity ($|T - T_\mathrm{crit}| < \SI{e-7}{K}$ deterministically crashed the
equation-of-state flash).

%=====================================================================
\section{Computational Efficiency}
\label{sec:efficiency}
%=====================================================================

All optimizations were verified against frozen 30-step solver-state dumps of a
W7-X conductor model; the accepted tolerance was floating-point summation-order noise
(worst observed relative deviation $2.6\times10^{-13}$; the tabulated-property step
$2.9\times10^{-8}$ by design, see below).

\subsection{Electric module}

\begin{itemize}
  \item Stiffness assembly via \code{scipy.sparse.bmat} block composition instead of
        per-block LIL writes (removed a \SI{15}{s} hotspot per 30 steps).
  \item A sparse \emph{reduction operator} $P$ (equipotential merging and
        fixed-potential elimination) is built once; each electric substep reduces the
        system as $P^{T}KP$ and $P^{T}b$ instead of densifying the full matrix.
  \item Electric \emph{topology caching}: coordinates, connectivity, incidence and
        contact matrices, conductance, inductance and boundary-condition structure are
        built once (unless the mesh adapts); only resistance and stiffness are rebuilt
        per substep.
  \item The sparse solve uses the default COLAMD column ordering; the previously
        forced natural ordering was $112\times$ slower on the reduced
        10\,160-DOF system.
  \item The number of electric substeps is derived from the physical time scales
        instead of being hard-coded.
\end{itemize}

\subsection{Thermal assembler vectorization}

The per-element Python loop over the mesh was eliminated entirely. All Gauss-point and
element matrix builders operate on whole-mesh batches of shape
$(n_\mathrm{elem}, n_\mathrm{dof}, n_\mathrm{dof})$; global assembly exploits the fact
that, for a fixed local (row, column) pair, destination band-storage columns are
distinct across elements, so each of the $12\times12$ local pairs is scattered with a
single collision-free vectorized addition. The known-term construction was rewritten as
a banded matrix--vector product over the 23 band diagonals. Function-call count per
30 steps dropped from 15.3 to 1.7 million.

\subsection{Linear algebra}

The Fortran-translated banded LU (\code{gredub}/\code{gbacsb}) was replaced by
LAPACK via \code{scipy.} \code{linalg.solve\_banded}, including a documented remapping between
the legacy band-storage layout and LAPACK layout (verified to $10^{-15}$ on captured
systems).

\subsection{Tabulated fluid properties}

CoolProp equation-of-state flashes dominated the remaining runtime. CoolProp's own
tabular backend proved unusable (unsupported in the installed high-level interface;
erroneous expansion/compressibility values through the low-level one). A custom
\code{\_FluidPropertyTable} was implemented instead: per property, a bicubic
\code{RectBivariateSpline} on a $\log T \times \log p$ grid
($400 \times 300$ points, \SIrange{2.4}{350}{K}, \SIrange{e5}{2.5e7}{Pa}), built once
per fluid ($\sim$\SI{1}{s}) from the reference equation of state. Accuracy versus the
full equation of state is $\le 10^{-6}$ at the 99th percentile; the table construction
additionally repairs a known NaN hole in the helium conductivity model by neighbor
diffusion. The property evaluation cost per 30 steps fell from 2.85 to \SI{0.53}{s}
($\approx 2.1\times$ overall solver speedup at that stage).

\subsection{Time integration: new steppers and adaptive error control}
\label{sec:timestepping}

The original code offered backward Euler ($\theta=1$) and Crank--Nicolson
($\theta=1/2$) with a degenerate ``adaptive'' controller whose error estimate diverged
at near-stagnant nodes, permanently pinning the step at its minimum. This layer was
redesigned:

\begin{itemize}
  \item \textbf{Galerkin method} (\code{METHOD=GAL}): the $\theta=2/3$ member of the
        $\theta$ family --- more accurate than backward Euler, less oscillatory than
        Crank--Nicolson.
  \item \textbf{Variable-step BDF2} (\code{METHOD=BDF2}): second-order backward
        differences with exact step-ratio coefficients
        \begin{equation*}
          \frac{M}{\Delta t_n}\Bigl(a_0 U^{n+1} - a_1 U^{n} + a_2 U^{n-1}\Bigr)
            + A\,U^{n+1} = b^{n+1},\qquad
          a_0 = \tfrac{1+2r}{1+r},\;
          a_1 = 1+r,\;
          a_2 = \tfrac{r^2}{1+r},
        \end{equation*}
        with $r = \Delta t_n / \Delta t_{n-1}$, a backward-Euler startup step and
        fully implicit source treatment. All one-step methods now carry a two-level
        solution history.
  \item \textbf{Adaptive error control} (\code{IADAPTIME=4}, tolerance
        \code{TOLTIME}): the local truncation error is estimated from the deviation of
        the new solution from its linear extrapolation in time (proportional to
        $\Delta t^2 \ddot U$), normalized \emph{per field} (each fluid's velocity,
        pressure, temperature; each solid's temperature) with a weighted
        root-mean-square norm and physical magnitude floors --- this is what makes the
        estimate non-degenerate. The controller applies
        $\Delta t \leftarrow 0.9\,\Delta t \sqrt{\varepsilon/\mathrm{err}}$ with growth
        clamps $[0.25, 1.5]$, honors \code{STPMIN}/\code{STPMAX}, and clamps the step to
        land exactly on requested output times and on square-wave heating-window edges
        (both would otherwise be skipped by a variable step). There is no step
        rejection (a state rollback is not possible in the present architecture), which
        the conservative safety factor and growth clamp compensate.
\end{itemize}

On the W7-X benchmark the controller behaves exactly as the physics demands: it holds
the \SI{1}{ms} floor during the thermal-hydraulic quenchback window ($t \lesssim
\SI{2}{s}$, where the error indicator sits an order of magnitude above tolerance),
opens to $\sim$\SI{35}{ms} by $t=\SI{4}{s}$ with the indicator riding at tolerance, and
saturates at \code{STPMAX} $=\SI{0.1}{s}$ in the hot-spot plateau.

\subsection{Aggregate results}

\begin{table}[ht]
  \centering
  \caption{Solver cost for 30 time steps of the W7-X case, 5\,081-node mesh,
           during the optimization campaign. Every stage verified bit-compatible
           against the previous one.}
  \label{tab:perf30}
  \begin{tabular}{lS[table-format=2.1]}
    \toprule
    Stage & {Wall time (\si{s})} \\
    \midrule
    Baseline (after runtime migration)                      & 61.7 \\
    Electric assembly, reduction operator, LAPACK banded    &  9.2 \\
    Banded known-term, vectorized assembly, COLAMD          & 22.1\textsuperscript{a} \\
    Full element-loop vectorization                          &  8.0 \\
    Tabulated fluid properties                               &  3.85 \\
    \bottomrule
  \end{tabular}

  \vspace{2pt}
  {\footnotesize \textsuperscript{a}~Intermediate stage measured on a finer
  profiling configuration; stages are cumulative in code but timings were taken as
  the campaign progressed.}
\end{table}

\begin{table}[ht]
  \centering
  \caption{End-to-end W7-X quench benchmark, \SI{25}{s} transient, 17\,000 linear
           elements (85\% refined in the $[60,100]$\,\si{m} zone), versus the
           THEA reference.}
  \label{tab:endtoend}
  \begin{tabular}{lccc}
    \toprule
     & Fixed $\Delta t = \SI{1}{ms}$ (BE) & Adaptive BDF2 (\code{TOLTIME}$=10^{-3}$) & THEA \\
    \midrule
    Time steps          & 25\,000            & 2\,102   & --- \\
    Wall time           & $\approx$\SI{4}{h} & \textbf{\SI{19}{min}} & $\approx$\SI{15}{min} \\
    Hot spot at \SI{2}{s}  & \SI{90.1}{K}    & \SI{89.6}{K}  & \SI{87.3}{K} \\
    Hot spot at \SI{10}{s} & ---              & \SI{140.9}{K} & \SI{138}{K} \\
    Peak hot spot          & ---              & \SI{140.9}{K} & $\approx$\SIrange{135}{140.5}{K}\textsuperscript{a} \\
    \bottomrule
  \end{tabular}

  \vspace{2pt}
  {\footnotesize \textsuperscript{a}~Spread of THEA's own convergence study across
  meshes and time-differencing methods; THEA-BDF2 with loose tolerance gives
  \SI{140.5}{K}.}
\end{table}

%=====================================================================
\section{Verification Methodology}
%=====================================================================

Because the refactoring proceeded on a mid-flight branch without a runnable regression
suite, verification gates were layered as the code matured:

\begin{enumerate}
  \item \textbf{Reproducible execution}: headless drivers for the TDD example cases
        and the W7-X benchmark (phase 0) allow every configuration to run
        non-interactively from the command line.
  \item \textbf{Static gate} (early phases): full-tree \code{py\_compile}, an import
        smoke test growing to 39 core modules, \code{pyflakes} on every new file, and
        dangling-reference greps for every moved or renamed symbol (including
        name-mangled \code{\_Class\_\_method} forms, which can be called cross-module).
  \item \textbf{Bit-compatibility gate} (performance phases): frozen 30-step solver
        state dumps compared element-wise after every optimization; acceptance
        threshold at floating-point noise ($\lesssim 10^{-12}$ relative).
  \item \textbf{Physical benchmark} (final phase): the W7-X quench case reproduces
        the THEA reference solution (hot-spot trajectory within
        \SIrange{1}{2}{K} at all checkpoints; self-limiting quench with the correct
        $\approx$\SI{140}{K} plateau). A backward-Euler regression of the new
        time-integration layer against the pre-change code was identical to
        $2\times10^{-13}$.
\end{enumerate}

%=====================================================================
\section{Outlook}
%=====================================================================

\begin{itemize}
  \item \textbf{Step rejection}: the adaptive controller currently cannot redo a step;
        adding a rollback of the conductor state would allow more aggressive growth
        factors and looser tolerances.
  \item \textbf{Remaining hotspots} (all now legitimate, none loop-bound):
        vectorized assembly, the electric transient solve, spline evaluation and the
        banded solve share the remaining per-step cost roughly equally.
  \item \textbf{Electric time integration}: the electric subproblem still uses the
        $\theta$ family only; \code{BDF2} there falls back to backward Euler.
\end{itemize}

\end{document}
